\subsection{What SFT can and cannot teach}

SFT is behavior cloning: it is appropriate when the desired action, format, tool syntax,
or reasoning style can be demonstrated. The dataset mixture therefore defines the target
policy. A system prompt repeated unchanged in every example is only a constant condition;
the data do not show how behavior should vary when that prompt changes. Likewise, one
demonstrated path receives probability even when several alternatives would be valid.

The more important limitation for the later comparison is distributional. During SFT,
every target follows the correct demonstrated prefix. During deployment, the model
conditions on its own earlier samples and may enter states absent from the demonstrations.
SFT can provide a strong initial policy but does not itself explore alternative solutions,
judge terminal success, or learn recovery from the model's own errors. RL addresses the
outcome question using sampled behavior; on-policy distillation instead supplies dense
teacher feedback on states encountered by the student.
